% Options for packages loaded elsewhere
\PassOptionsToPackage{unicode}{hyperref}
\PassOptionsToPackage{hyphens}{url}
\PassOptionsToPackage{dvipsnames,svgnames*,x11names*}{xcolor}
%
\documentclass[
  11pt,
]{article}
\usepackage{lmodern}
\usepackage{amssymb,amsmath}
\usepackage{ifxetex,ifluatex}
\ifnum 0\ifxetex 1\fi\ifluatex 1\fi=0 % if pdftex
  \usepackage[T1]{fontenc}
  \usepackage[utf8]{inputenc}
  \usepackage{textcomp} % provide euro and other symbols
\else % if luatex or xetex
  \usepackage{unicode-math}
  \defaultfontfeatures{Scale=MatchLowercase}
  \defaultfontfeatures[\rmfamily]{Ligatures=TeX,Scale=1}
\fi
% Use upquote if available, for straight quotes in verbatim environments
\IfFileExists{upquote.sty}{\usepackage{upquote}}{}
\IfFileExists{microtype.sty}{% use microtype if available
  \usepackage[]{microtype}
  \UseMicrotypeSet[protrusion]{basicmath} % disable protrusion for tt fonts
}{}
\makeatletter
\@ifundefined{KOMAClassName}{% if non-KOMA class
  \IfFileExists{parskip.sty}{%
    \usepackage{parskip}
  }{% else
    \setlength{\parindent}{0pt}
    \setlength{\parskip}{6pt plus 2pt minus 1pt}}
}{% if KOMA class
  \KOMAoptions{parskip=half}}
\makeatother
\usepackage{xcolor}
\IfFileExists{xurl.sty}{\usepackage{xurl}}{} % add URL line breaks if available
\IfFileExists{bookmark.sty}{\usepackage{bookmark}}{\usepackage{hyperref}}
\hypersetup{
  pdftitle={Literature Review: Solution of Linear ODEs with Delta Function as Load},
  pdfauthor={Denis Pleshkov},
  colorlinks=true,
  linkcolor=Maroon,
  filecolor=Maroon,
  citecolor=Blue,
  urlcolor=Blue,
  pdfcreator={LaTeX via pandoc}}
\urlstyle{same} % disable monospaced font for URLs
\usepackage[margin=1in]{geometry}
\usepackage{longtable,booktabs}
% Correct order of tables after \paragraph or \subparagraph
\usepackage{etoolbox}
\makeatletter
\patchcmd\longtable{\par}{\if@noskipsec\mbox{}\fi\par}{}{}
\makeatother
% Allow footnotes in longtable head/foot
\IfFileExists{footnotehyper.sty}{\usepackage{footnotehyper}}{\usepackage{footnote}}
\makesavenoteenv{longtable}
\setlength{\emergencystretch}{3em} % prevent overfull lines
\providecommand{\tightlist}{%
  \setlength{\itemsep}{0pt}\setlength{\parskip}{0pt}}
\setcounter{secnumdepth}{-\maxdimen} % remove section numbering

\title{Literature Review: Solution of Linear ODEs with Delta Function as
Load}
\author{Denis Pleshkov (\href{mailto:std.approach@gmail.com}{std.approach@gmail.com})}
\date{August 2026}

\begin{document}
\maketitle

\hypertarget{abstract}{%
\subsubsection{Abstract}\label{abstract}}

This literature review examines how linear time-invariant (LTI) ordinary
differential equations (ODEs) with impulsive forcing (represented by the
Dirac delta function and its derivatives) are treated across
differential equations, vibration theory, and control theory. Motivated
by the practical need to find proper solution approaches and closed-form
formulas for such systems, we surveyed 100+ sources. Critical gaps in
the literature are identified --- limited treatment of derivatives of
delta, and lack of unified computational frameworks.

\hypertarget{keywords}{%
\subsubsection{Keywords}\label{keywords}}

Literature review, delta function, differential equation, ODE, impulse
response

\hypertarget{introduction}{%
\subsubsection{Introduction}\label{introduction}}

The dynamics of evolving processes are often subject to abrupt changes,
such as: impact of a hammer on a beam, a bat striking a ball, a bolt of
lightning striking a tower.

Such short-term perturbations are frequently treated as instantaneous
events, often modeled as ``impulses'' (see also Cohen, p.~13, for an
analogous billiard-ball example). Logan (p.166) ``Many physical and
biological processes have source terms that act at a single instant of
time. For example, we can idealize an injection of medicine (a''shot``)
into the blood stream as occurring at a single instant; a mechanical
system, for example, a damped spring--mass system in a shock absorber on
a car can be given an impulsive force by hitting a bump in the road; the
switch in an electrical circuit can be closed only for an instant, which
leads to an impulsive, applied voltage''. According to Rao (p.~381)
``The simplest form is the impulsive force a force that has a large
magnitude F and acts for a very short time''. The system's response to
such a force is termed the impulse response function (IRF).
Mathematically, an impulse can be represented within an initial value
problem (IVP) by incorporating the Dirac delta function as the external
forcing term. An IVP consists of an ODE together with the system's state
at some initial time; its solution is the unique function satisfying
both. The impulse response of a system is defined as its output in
response to an input \(\delta(t)\), assuming the system is initially at
rest.

We are searching the literature for solutions to LTI ODEs with a
discontinuous right-hand side, including the delta function and its
derivatives. The Dirac delta function is a well-known generalized
function (distribution) used to model impulsive phenomena. Its
properties are discussed extensively in the literature, including Arfken
(pp.76-79), Bottega (p.\,233), Chasnov (p.\,58), Nagy (p.\,196-201),
Weber (p.\,86-90), Zill (p.\,328-330).

While seeking a general method to solve such systems, we found that
existing literature primarily offers solutions for specific first- and
second-order ODEs and almost no books present the common (closed form)
solution.

The sources reviewed here were identified through a manual survey of
over 100 textbooks and articles in differential equations, vibration
theory, and control theory, selected for their treatment of impulsive or
discontinuous forcing in linear time-invariant systems. Section 1
establishes the equivalence pattern observed across this literature;
Section 2 classifies the surveyed sources into five categories based on
how explicitly they state this equivalence; Section 3 draws conclusions
and identifies the resulting gaps.

Beyond the sources analyzed in detail below, a broader survey of the
Control Theory literature identified more than fifty additional
textbooks that discuss the Impulse Response Function (IRF) --- the
solution of an LTI ODE with zero Initial Condition (IC) and an impulse
delta function as load --- without engaging the equivalence question
directly. This broader bibliography is provided in the Supplementary
Bibliography following the References, for readers seeking a starting
point for further study of Control Theory.

\hypertarget{equivalence-through-initial-conditions-modification}{%
\subsubsection{1 Equivalence through initial condition's
modification}\label{equivalence-through-initial-conditions-modification}}

Several textbooks provide analytical solutions for first/second order
LTI ODEs with a Dirac delta forcing function. Examples include Nagy
(pp.~203-209: ``The Impulse Response Function''), Ogata (p.\,163:
``Unit-Impulse Response of First-Order Systems''; p.178: ``Impulse
Response of Second-Order Systems''), Rao (p.\,382: ``4.5.1 Response to
an Impulse''), and Zill (p.\,330: ``Two Initial-Value Problems'').

Baruh (p.9): ``Impulsive forces cause sudden changes in velocity with
very little (or negligible) change in position''

Other authors explicitly note that the solution of an IVP with a delta
load coincides with the solution of the corresponding homogeneous ODE
subject to modified initial conditions (IC). A brief survey of such
observations follows.

Rao (p.\,407: ``Unit Impulse Response of a First-Order System'') states
the equivalence for 1st order system

\[
\begin{cases}
y' + a y = F \delta(t), \\
y(0) = 0
\end{cases}
\;\equiv\;
\begin{cases}
y' + a y = 0, \\
y(0) = F
\end{cases}
\]

Weber (p.\,733: ``Impulsive Force'') notes the analogous result

\[
\begin{cases}
m x'' = P \delta(t), \\
x(0) = 0, \\
x'(0) = 0
\end{cases}
\;\equiv\;
\begin{cases}
m x'' = 0, \\
x(0) = 0, \\
x'(0) = P/m
\end{cases}
\]

Balachandran (p.~301: ``Similarity to Response to Initial Velocity''),
Bottega (pp.\,235--236: ``problem of interest becomes equivalent to the
problem of free vibrations with the initial conditions''), Genta
(p.\,179-180: ``The position x0 after the impulse is equal to that
before the impulse, while the velocity v0 is equal to the one before the
impulse plus an increment''), Meirovitch (pp.\,160--161: ``we conclude
that the effect of a unit impulse at t = 0 is to produce an equivalent
initial velocity''), Schiff (p.\,83: ``indicating the instantaneous jump
in velocity at t=0, from a rest state to the value v0''), all remarked
that

\[
\begin{cases}
m x'' + c x' + k x = f_0 \delta(t), \\
x(0) = 0, \\
x'(0) = 0
\end{cases}
\;\equiv\;
\begin{cases}
m x'' + c x' + k x = 0, \\
x(0) = 0, \\
x'(0) = f_0/m
\end{cases}
\]

Chasnov (p.\,61) provides a formula for changing the IC of a
second-order LTI ODE. The LTI IVP with discontinuous right side can also
be viewed as a special case of an impulsive differential equation (see
Benchohra, Henderson, and Ntouyas).

The above examples suggest a general pattern (though a formal proof
appears to be missing in the literature): an IVP forced by a delta
function may be equivalent to a homogeneous IVP with a shifted IC.

\hypertarget{detailed-literature-classification}{%
\subsubsection{2 Detailed literature
classification}\label{detailed-literature-classification}}

We surveyed the following sources and classified each by how explicitly
it addresses the equivalence between an impulse-forced IVP and a
homogeneous IVP with modified initial conditions.

\begin{longtable}[]{@{}llll@{}}
\toprule
\begin{minipage}[b]{0.22\columnwidth}\raggedright
Category\strut
\end{minipage} & \begin{minipage}[b]{0.22\columnwidth}\raggedright
Criterion\strut
\end{minipage} & \begin{minipage}[b]{0.22\columnwidth}\raggedright
Count\strut
\end{minipage} & \begin{minipage}[b]{0.22\columnwidth}\raggedright
Representative authors\strut
\end{minipage}\tabularnewline
\midrule
\endhead
\begin{minipage}[t]{0.22\columnwidth}\raggedright
1\strut
\end{minipage} & \begin{minipage}[t]{0.22\columnwidth}\raggedright
Solves the impulse-forced system without stating the equivalence\strut
\end{minipage} & \begin{minipage}[t]{0.22\columnwidth}\raggedright
16\strut
\end{minipage} & \begin{minipage}[t]{0.22\columnwidth}\raggedright
Alam, Boyce, Kreyszig, Ricardo\strut
\end{minipage}\tabularnewline
\begin{minipage}[t]{0.22\columnwidth}\raggedright
2\strut
\end{minipage} & \begin{minipage}[t]{0.22\columnwidth}\raggedright
Mentions the change of IC qualitatively, without a formula\strut
\end{minipage} & \begin{minipage}[t]{0.22\columnwidth}\raggedright
13\strut
\end{minipage} & \begin{minipage}[t]{0.22\columnwidth}\raggedright
Antsaklis, Brogliato, Logan, Shabana\strut
\end{minipage}\tabularnewline
\begin{minipage}[t]{0.22\columnwidth}\raggedright
3\strut
\end{minipage} & \begin{minipage}[t]{0.22\columnwidth}\raggedright
Gives an explicit formula for specific equation orders\strut
\end{minipage} & \begin{minipage}[t]{0.22\columnwidth}\raggedright
30\strut
\end{minipage} & \begin{minipage}[t]{0.22\columnwidth}\raggedright
Angeles, Chopra, Meirovitch, Trench\strut
\end{minipage}\tabularnewline
\begin{minipage}[t]{0.22\columnwidth}\raggedright
4\strut
\end{minipage} & \begin{minipage}[t]{0.22\columnwidth}\raggedright
Gives a formula for a general n-th order equation, delta-only load\strut
\end{minipage} & \begin{minipage}[t]{0.22\columnwidth}\raggedright
5\strut
\end{minipage} & \begin{minipage}[t]{0.22\columnwidth}\raggedright
Adkins, Beneš, Camporesi, Filippov\strut
\end{minipage}\tabularnewline
\begin{minipage}[t]{0.22\columnwidth}\raggedright
5\strut
\end{minipage} & \begin{minipage}[t]{0.22\columnwidth}\raggedright
Closed-form solution for a sum of derivatives of delta\strut
\end{minipage} & \begin{minipage}[t]{0.22\columnwidth}\raggedright
0\strut
\end{minipage} & \begin{minipage}[t]{0.22\columnwidth}\raggedright
--- (open gap)\strut
\end{minipage}\tabularnewline
\bottomrule
\end{longtable}

\hypertarget{solution-without-mentioning-the-equivalence-i.e.-unit-impulse-is-equal-to-change-ic}{%
\paragraph{2.1 Solution without mentioning the equivalence (i.e.~unit
impulse is equal to change
IC)}\label{solution-without-mentioning-the-equivalence-i.e.-unit-impulse-is-equal-to-change-ic}}

Alam (p.\,218 1st order, pp.\,270-274 2nd order), Asadi (pp.~62-73),
Benaroya (p.\,147 IRF 2nd order), Boyce (p.272), Campbell (p.263), Dorf
(p.327), Etkin (pp.\,73-76), Goode (pp.\,708--710), Gupta (pp.\,72, 86),
Holmes (p.179), Jack (p.\,575), Jeffrey (p.412), Korman (p.\,160),
Kreyszig (pp.\,227--230), Neil (pp.102-104), Ricardo (pp.215-216).

\hypertarget{mention-the-change-of-initial-condition-without-giving-an-explicit-formula}{%
\paragraph{2.2 Mention the change of initial condition without giving an
explicit
formula}\label{mention-the-change-of-initial-condition-without-giving-an-explicit-formula}}

Anderson (p.~8: momentum argument for 2nd order), Antsaklis (p.~67:
``the impulse response of a linear, time-invariant, continuous time
system with integral representation is equal to the kernel of the
integral representation of the system''), Benaroya (p.19 due to
``principle of conservation of linear momentum'' governs the change the
velocity via impulse `During collisions large forces act resulting in
almost instantaneous changes in velocity and therefore in linear
momentum'), Brogliato (p.2: ``impulsive forces imply a discontinuity in
the velocity while positions remain continuous''; p.7: in mechanical
systems, continuous positions and discontinuous velocities are produced
by impulsive forces, and vice versa``), Howell (p.~560:''Observe that
using a delta function force leads to the velocity changing instantly
from one constant to another``), James (p.345, p.365), Kausel (p.~82:''
the impulse imparted on a mass m abruptly changes its velocity from zero
to V = 1/m``), Logan (pp.169-170, p.172), MacCluer (p.374:''Notice that
the solution in equation \ldots{} doesn't actually satisfy the initial
condition``), McOwen (pp.~98-99), Peterson (p.~363:''discontinuous
forcing function causes a jump in the velocity of the mass``), Ram
(p.~22-5), Shabana (p.~40:''This result indicates that the effect of the
impulsive force, which acts over a very short time duration on a system
which is initially at rest, can be accounted for by considering the
motion of the system with initial velocity 11m and zero initial
displacement.")

\hypertarget{provide-an-explicit-formula-for-changing-initial-conditions-for-specific-equations}{%
\paragraph{2.3 Provide an explicit formula for changing initial
conditions for specific
equations}\label{provide-an-explicit-formula-for-changing-initial-conditions-for-specific-equations}}

Angeles (p.115: ``consequence, the ball undergoes a finite change in its
velocity'', p.116: change IC for IRF for 1st order system, p.119: change
IC for IRF for 2nd order system, pp.132/136: derivative of delta as
load), Balachandran (p.~301: IRF for 2nd order system), Baruh (p.259:
change IC for 2nd order system due to impulse load), Beards (p.66: ``the
impulse F, acting on a body will result in a sudden change in its
velocity without an appreciable change in its displacement. Thus the
motion of a single degree of freedom system excited by an impulse F,
corresponds to free vibration with initial conditions x = 0 and v0, =
F/m at t = 0''), Chopra (Demonstrates that discontinuous forcing via
impulse is mathematically equivalent to modified initial conditions;
p.121 exact formulae for change IC for linear oscillator (from zero IC);
p.616 example of impulse response for MDOF), Cooper (Provides explicit
formula for piecewise smooth derivatives showing jump
\(\leftrightarrow\) delta connection. Foundational reference proving
mathematically that impulses (delta forces) arise naturally from
discontinuities; p.5 change IC for pendulum was in rest), Duffy (p.93
``This avoids the problem of the Green's function not satisfying all of
the initial conditions.'' + (3.1.7) provide the initial condition
delivers the IRF as free motion., p.166, p.284), Edwards (p.~500, 2nd
order with time shift), Esfandiari (p.~57: change IC for IRF for 2nd
order system/``when impulsive forces are present in the system, initial
values and initial conditions are indeed different'', p.343: ``Impulse
Response of First-Order Systems'', p.351/359: ``Impulse Response of
Second-Order Systems''), Fairman (page 31, In addition we see from
(1.81) that the zero-input response equals the impulse response when the
initial state is x(0) = B (IRF is equal to some non-zero IC)), Franklin
(p.110: change IC on 1st order system), Haddad (p.50 ``we can always
reproduce the impulsive response with the free response by setting x(0)
= Bv''), Hallauer (p.~158 change IC for 1st order system), Inman p.219:
``impulsive load for the system initially at rest is calculated by
recalling from physics that an impulse imparts a change in momentum to a
body'', p.557: IRF for PDE/string), Iyengar (p.~87: change IC for SDOF
whilst delta load, p.121 IRF for MDOF), Jazar (pp.~173, p.188: ``Impulse
will only change the initial conditions, and hence, the response of a
multi-DOF system will be the transient response to a new set of initial
conditions''), Kabe (pp.149-150: ``it would appear that applying a unit
impulse at t = 0 is equivalent to giving the system an initial
velocity'', p.163), Klee (p.170: change IC for 1st order system), Lathi
(pp.164-165: ``Find the impulse response h(t) for a system specified by
(D\^{}2 +5D+6)y(t) = (D+1)x(t)''), Luintel (p.215: ``velocity of the
system immediately after the application of impulse I is I/m'', 507),
Meirovitch (p.161: ``effect of a unit impulse at t = 0 is to produce an
equivalent initial velocity''), Nielsen (p.24: ``Consequently, an
impulsive load causes a discontinuous change of the velocity'', p.63:
solution for IRF for MDOF system), Palm (p.96: change IC for 1st/2nd
order system, p.116: ``Impulse Response of Second-Order Models''),
Polking (p.231: IRF for 2nd order system), Schiff (p.82-83 2nd order
system with impulse is equal to free vibration of the same system but
changed IC), Silva (p.87 for 2nd order system for IRF mentioned change
IC, ``The Riddle of Zero Initial Conditions''), Sinha (p.69 mentioned
changing IC for solution of 1st order ODE with unit impulse as load),
Thorby (p.50: ``unit impulse of force is applied to it \ldots{} that if
dv is the change in velocity, then dv=1/m, but the change in isplacement
is negligible''), Trench (pp.~478--480: solution some 2nd order system
with impulse load), Tse (p.~54: change IC for 2nd order system with
delta load)

\hypertarget{supply-a-formula-for-the-case-where-the-load-contains-only-the-delta-function-or-its-derivatives-for-general-n}{%
\paragraph{\texorpdfstring{2.4 Supply a formula for the case where the
load contains only the delta function (or its derivatives) for general
\(n\)}{2.4 Supply a formula for the case where the load contains only the delta function (or its derivatives) for general n}}\label{supply-a-formula-for-the-case-where-the-load-contains-only-the-delta-function-or-its-derivatives-for-general-n}}

\begin{itemize}
\tightlist
\item
  Adkins (p.\,319: provide formulae for changing IC to calculate IRF as
  free response for LTI ODE n-th order)
\item
  Beneš (article; for LTI ODE with delta load, provided the whole
  formula for changing IC - (9); the authors were very close to provide
  the whole formula for LTI ODE with the delta and its derivatives as
  load)
\item
  Camporesi (p.2: for n-th order equation in the form giving a formulae
  for changing IC to get IRF as free vibration, p.5: example for 1st
  order system)
\item
  Campos (p.~166: change IC for oscillator with first derivative of
  delta, p.167: exact solution for oscillator with odd/even derivative
  of delta as load)
\item
  Filippov (the authors used the approach from Beneš; p.20: author
  provided the trick how to convert original ODE with delta and its
  derivatives as right side into system of 1st order ODE with changed
  right side that contains only regular function, not general ones. BUT
  the conversion formula missed the important part ``how to change IC to
  solve the changed system''. We need to conclude that Filippov's result
  is not usable for engineering tasks)
\end{itemize}

\hypertarget{closed-form-solution-for-an-lti-ode-with-a-sum-of-derivatives-of-the-dirac-delta-as-the-forcing-function}{%
\paragraph{2.5 Closed-form solution for an LTI ODE with a sum of
derivatives of the Dirac delta as the forcing
function}\label{closed-form-solution-for-an-lti-ode-with-a-sum-of-derivatives-of-the-dirac-delta-as-the-forcing-function}}

To the best of our knowledge, no such solution exists in the literature
(see Conclusion).

\hypertarget{conclusion}{%
\subsubsection{3 Conclusion}\label{conclusion}}

Across the sources surveyed, a consistent pattern emerges: the solution
of an LTI ODE driven by a Dirac delta forcing term coincides with the
solution of the corresponding homogeneous equation under a shifted
initial condition (Section 1). The literature treats this equivalence
unevenly, however. As the classification in Section 2 shows, most
sources either use the equivalence implicitly, without stating it
(category 1), or mention it qualitatively without an explicit formula
(category 2); a smaller subset derive an explicit formula for a specific
order of equation (category 3); and fewer still address the general case
of an n-th order equation forced by the delta function alone (category
4). None of the surveyed sources provide a closed-form solution for the
most general case --- an LTI ODE forced by a sum of derivatives of the
Dirac delta (category 5).

This gap points to two open needs in the literature: a treatment of
derivatives of the delta function as forcing terms that is as thorough
as the treatment of the delta function itself, and a unified
computational framework that covers the full range of derivative and
equation orders within a single formula, rather than the
equation-by-equation formulas found in category 3.

\hypertarget{references}{%
\subsubsection{REFERENCES}\label{references}}

Adkins, W. A., \& Davidson, M. G. (2012). \emph{Ordinary differential
equations}. Springer Science+Business Media.
https://doi.org/10.1007/978-1-4614-3618-8

Alam, J., Hu, G., Babu, H. M. H., \& Xu, H. (2023). \emph{Control
engineering: Theory and applications}. CRC Press.
https://doi.org/10.1201/9781003293859

Anderson, B., \& Rufer, S. (2018, August 13). \emph{Control theory: A
brief introduction}. Bruin Racing, Baja SAE, University of California,
Los Angeles. https://doi.org/10.13140/RG.2.2.14805.17129

Angeles, J. (2011). \emph{Dynamic response of linear mechanical systems:
Modeling, analysis and simulation}. Springer Science+Business Media.
https://doi.org/10.1007/978-1-4419-1027-1

Antsaklis, Panos J.; Michel, Anthony N. \emph{A Linear Systems Primer}.
Birkhäuser, 2007. ISBN: 9780817644604

Arfken, G. B., Weber, H. J., \& Harris, F. E. (2011). Mathematical
methods for physicists: A comprehensive guide (7th ed.). Academic Press
/ Elsevier. Print ISBN: 978-0-12-384654-9

Asadi, F., Bolanos, R. E., \& Rodríguez, J. (2019). \emph{Feedback
control systems: The MATLAB®/Simulink® approach} (Synthesis Lectures on
Control and Mechatronics, Lecture \#5). Morgan \& Claypool Publishers.
https://doi.org/10.2200/S00909ED1V01Y201903CRM005

Balachandran, B., \& Magrab, E. B. (2019). \emph{Vibrations} (3rd ed.).
Cambridge University Press. https://doi.org/10.1017/9781108615839

Baruh, H. (2015). \emph{Applied dynamics}. CRC Press, Taylor \& Francis
Group. (ISBN: 978-1-4822-0734-7) https://doi.org/10.1201/b18272

Beards, C. F. (1996). \emph{Structural vibration: Analysis and damping}.
Arnold; Halsted Press. (ISBN: 0340645806, 9780340645802, 0470235861,
9780470235867)

Benaroya, H., Nagurka, M., \& Han, S. (2017). \emph{Mechanical
vibration: Analysis, uncertainties, and control} (4th ed.). CRC Press,
Taylor \& Francis Group. (ISBN: 978-1-4987-5265-7)
https://doi.org/10.1201/b22347

Benchohra, Mouffak; Henderson, Johnny; Ntouyas, Sotiris K.
\emph{Impulsive Differential Equations and Inclusions}. Hindawi
Publishing, 2006.

Beneš, K. (1978). On modelling dynamic systems excited by the Dirac
function. \emph{Sborník prací Přírodovědecké fakulty University
Palackého v Olomouci. Matematika, 17}(1), 123--129.
http://dml.cz/dmlcz/120062

Bottega, W. J. (2006). \emph{Engineering vibrations}. CRC Press. (ISBN:
9780849334207, 0849334209)

Boyce, W. E., \& DiPrima, R. C. (2017). \emph{Elementary differential
equations and boundary value problems} (11th ed.). John Wiley \& Sons,
Inc.~(ISBN: 978-1-119-37792-4)

Brogliato, Bernard. \emph{Nonsmooth Mechanics: Models, Dynamics and
Control} (3rd ed.). Springer, 2015.

Campbell, S. L., \& Haberman, R. (2008). Introduction to differential
equations with dynamical systems. Princeton University Press. ISBN:
978-0-691-12474-2 (hardcover)

Camporesi, Carlo. \emph{An Introduction to Linear Ordinary Differential
Equations Using the Impulsive Response Method and Factorization}. 2019.

Campos, L. M. B. C. (2020). \emph{Linear differential equations and
oscillators} (Vol. 4). CRC Press, Taylor \& Francis Group. (ISBN:
978-0-367-13718-2) https://doi.org/10.1201/9780429028984

Chasnov, J. R. (2009--2016). \emph{Introduction to differential
equations: Lecture notes for MATH 2351/2352}. The Hong Kong University
of Science and Technology.

Chopra, A. K. (2020). \emph{Dynamics of structures: Theory and
applications to earthquake engineering} (5th ed., SI units). Pearson
Education Limited. (ISBN: 978-1-292-24918-6)

Cohen, A. M. (2007). \emph{Numerical methods for Laplace transform
inversion}. Springer Science+Business Media. (ISBN: 9780387282619,
0387282610) https://doi.org/10.1007/978-0-387-68855-8

Cooper, David. \emph{Distribution Theory}. 2000.

Dorf, R. C., \& Bishop, R. H. (2008). \emph{Modern control systems:
Solution manual} (11th ed.). Pearson Education, Inc.~(ISBN:
0-13-227029-3)

Duffy, D. G. (2015). \emph{Green's functions with applications} (2nd
ed.). CRC Press, Taylor \& Francis Group. (ISBN: 978-1-4822-5103-6)
https://doi.org/10.1201/b17973

Edwards, C. H., Penney, D. E., \& Calvis, D. (2016). \emph{Differential
equations and boundary value problems: Computing and modeling} (5th ed.,
Global ed.). Pearson Education Limited. (ISBN: 978-1-292-10877-3)

Esfandiari, R. S., \& Lu, B. (2014). \emph{Modeling and analysis of
dynamic systems} (2nd ed.). CRC Press, Taylor \& Francis Group. (ISBN:
978-1-4665-7495-3) https://doi.org/10.1201/b16443

Etkin, B. (2005). \emph{Dynamics of atmospheric flight}. Dover
Publications, Inc.~(ISBN: 0-486-44522-4) (Original work published 1972)

Fairman, Frederick W. \emph{Linear Control Theory: The State Space
Approach}. John Wiley \& Sons, 1998. ISBN: 0-471-97489-7

Filippov, A. F. (1988). \emph{Differential equations with discontinuous
righthand sides} (F. M. Arscott, Ed.). Springer-Science+Business Media,
B.V. (ISBN: 978-90-481-8449-1) https://doi.org/10.1007/978-94-015-7793-9
(Original work published 1988)

Franklin, G. F., Powell, J. D., \& Emami-Naeini, A. (2015).
\emph{Feedback control of dynamic systems} (7th ed., Global ed.).
Pearson Education Limited. (ISBN: 978-1-292-06890-9)

Genta, G. (2009). \emph{Vibration dynamics and control}. Springer
Science+Business Media, LLC. (ISBN: 978-0-387-79579-9, 9780387795805)
https://doi.org/10.1007/978-0-387-79580-5

Goode, S. W., \& Annin, S. A. (2015). \emph{Differential equations and
linear algebra} (4th ed.). Pearson Education, Inc.~(ISBN:
978-0-321-96467-0)

Gupta, A., \& Verma, Y. P. (2020). \emph{Automatic control engineering}.
I.K. International Pvt. Ltd.~(ISBN: 978-93-89583-74-8)

Haddad, Wassim M.; Chellaboina, Vijaysekhar; Hui, Qing.
\emph{Nonnegative and Compartmental Dynamical Systems}. Oxford
University Press, 2009. ISBN: 978-0-691-14411-5

Hallauer, William L. \emph{Linear Time-Invariant Dynamic Systems}. John
Wiley \& Sons, 2016.

Holmes, M. H. (2023). Introduction to differential equations (3rd ed.).
XanEdu. ISBN: 978-1-71147-191-4

Howell, K. B. (2020). \emph{Ordinary differential equations: An
introduction to the fundamentals} (2nd ed.). CRC Press, Taylor \&
Francis Group. (ISBN: 978-1-138-60583-1)
https://doi.org/10.1201/9780429347429

Inman, Daniel J. \emph{Engineering Vibration} (4th ed.). Pearson
Education, Inc., 2014. ISBN: 978-0-13-287169-3

Iyengar, R. N. (2019). \emph{Elements of mechanical vibration}. I.K.
International Pvt. Ltd.~(ISBN: 978-93-89633-34-4)

Jack, H. (2015). \emph{Dynamic system modeling and control}. Hugh Jack.
(ISBN: 978-1-5089-9525-8)

James, G., Dyke, P., Burley, D., Clements, D., Craven, M., Reis, T.,
Searl, J., Stander, J., Steele, N., \& Wright, J. (2018). \emph{Advanced
modern engineering mathematics} (5th ed.). Pearson Education Limited.
(ISBN: 978-1-292-17434-1)

Jazar, R. N., \& Marzbani, H. (2024). \emph{Vehicle vibrations: Linear
and nonlinear analysis, optimization, and design}. Springer Nature
Switzerland AG. (ISBN: 978-3-031-43485-3)
https://doi.org/10.1007/978-3-031-43486-0

Jeffrey, A. (2002). Advanced engineering mathematics. Harcourt/Academic
Press. ISBN-10: 0-12-382592-X

Kabe, A. M., \& Sako, B. H. (2020). \emph{Structural dynamics:
Fundamentals and advanced applications} (Vol. 1). Academic Press, an
imprint of Elsevier. (ISBN: 978-0-12-821614-9)
https://doi.org/10.1016/C2019-0-00137-8

Kausel, Eduardo. \emph{Advanced Structural Dynamics}. Cambridge
University Press, 2017. ISBN: 978-1-107-17151-0.
https://doi.org/10.1017/9781316761403

Klee, H., \& Allen, R. (2011). \emph{Simulation of dynamic systems with
MATLAB® and Simulink®} (2nd ed.). CRC Press, Taylor \& Francis Group.
(ISBN: 978-1-4398-3674-3) https://doi.org/10.1201/b10495

Korman, P. L. (2019). \emph{Lectures on differential equations}. MAA
Press, an imprint of the American Mathematical Society. (ISBN:
978-1-4704-5173-8)

Kreyszig, E. (2011). \emph{Advanced engineering mathematics} (10th ed.).
John Wiley \& Sons, Inc.~(ISBN: 978-0-470-45836-5)

Lathi, B. P., \& Green, R. A. (2018). \emph{Linear systems and signals}
(3rd ed.). Oxford University Press. (ISBN: 978-0-19-020017-6)

Logan, J. D. (2015). \emph{A first course in differential equations}
(3rd ed.). Springer-Verlag. (ISBN: 978-3-319-17851-6)
https://doi.org/10.1007/978-3-319-17852-3

Luintel, M. C. (2024). \emph{Textbook of mechanical vibrations}.
Springer Nature Singapore Pte Ltd.~(ISBN: 978-981-99-3613-7)
https://doi.org/10.1007/978-981-99-3614-4

McOwen, R. (2012). Worldwide differential equations with linear algebra
(1st ed.). Worldwide Center of Mathematics, LLC. ISBN-10: 0-9842071-2-0

Meirovitch, L. (2001). \emph{Fundamentals of vibrations} (International
ed.). McGraw-Hill. (ISBN: 0-07-118174-1)

Nagy, G. (n.d.). \emph{Ordinary differential equations}. Mathematics
Department, Michigan State University.

Nielsen, S. R. K. (2004). \emph{Vibration theory, Vol. 1: Linear
vibration theory} (3rd ed.). Department of Civil Engineering, Aalborg
University. (U/ Vol. U2004-1)

O'Neil, Peter V. Advanced Engineering Mathematics, SI. SI ed., 8th ed.,
Cengage Learning, 2018. ISBN-13: 978-1-337-27452-4

Ogata, K. (2010). \emph{Modern control engineering} (5th ed.). Pearson
Education, Inc.~(ISBN-13: 978-0-13-615673-4)

Palm, W. J., III. (2010). \emph{System dynamics} (2nd ed.). McGraw-Hill.
(ISBN: 978-0-07-352927-1)

Peterson, G. L., \& Sochacki, J. S. (2014). \emph{Linear algebra \&
differential equations} (Pearson New International ed.). Pearson
Education Limited. (ISBN: 978-1-269-37450-7)

Polking, J., Boggess, A., \& Arnold, D. (2006). \emph{Differential
equations with boundary value problems} (2nd ed.). Pearson Prentice
Hall. (ISBN: 0-13-186236-7)

Ram, B. (2009). \emph{Engineering mathematics}. Pearson Education.
(ISBN: 978-81-317-2691-4)

Rao, S. S. (2011). \emph{Mechanical vibrations} (5th ed.). Pearson
Education, Inc.~(ISBN-13: 978-0-13-212819-3)

Ricardo, H. J. (2020). A modern introduction to differential equations
(3rd ed.). Academic Press. https://doi.org/10.1016/C2018-0-02231-8 Print
ISBN-13: 978-0-12-823417-4

Schiff, Joel L. (1999). \emph{The Laplace transform: Theory and
applications}. Springer-Verlag New York, Inc.~(ISBN: 0-387-98698-7)
https://doi.org/10.1007/978-0-387-22757-3

Shabana, A. A. (1997). Vibration of discrete and continuous systems (2nd
ed.). Springer-Verlag. https://doi.org/10.1007/978-1-4612-4036-5 Print
ISBN-13: 978-1-4612-8474-1

Sinha, N.K., \& Ananthkrishnan, N. (2022). \emph{Elementary flight
dynamics with an introduction to bifurcation and continuation methods}
(2nd ed.). CRC Press. https://doi.org/10.1201/9781003096801

Thorby, D. (2008). \emph{Structural dynamics and vibration in practice:
An engineering handbook}. Butterworth-Heinemann, an imprint of Elsevier.
(ISBN: 978-0-7506-8002-8)
https://doi.org/10.1016/B978-0-7506-8002-8.X0001-1

Trench, W. F. (2024). \emph{Elementary differential equations with
boundary values problems}. LibreTexts. Retrieved December 19, 2024, from
https://LibreTexts.org

Tse, F. S., Morse, I. E., \& Hinkle, R. T. (2018). \emph{Mechanical
vibrations: Theory and applications} (2nd ed.). CBS Publishers \&
Distributors Pvt. Ltd.~(eISBN: 978-93-879-6458-7)

Weber, H. J., \& Arfken, G. B. (2003). \emph{Essential mathematical
methods for physicists}. Academic Press. (ISBN: 978-0-12-059877-9)
https://doi.org/10.1016/B978-0-12-059877-9.X5000-7

Zill, D. G. (2023). A first course in differential equations with
modeling applications (12th ed.). Cengage Learning. Hardcover ISBN:
978-0-357-76019-2

\hypertarget{supplementary-bibliography}{%
\subsubsection{Supplementary
Bibliography}\label{supplementary-bibliography}}

Many additional books related to Control Theory discuss the IRF, the
solution of an LTI ODE with zero IC and an impulse delta function as
load, without directly addressing the equivalence question examined in
this review. Here is a short list of such books, which could be used as
a starting point for further studying Control Theory:

\begin{itemize}
\tightlist
\item
  Abdallah - Feedback Control Systems MATLAB Simulink Approach
\item
  Alam, Jahangir - Control Engineering Theory and Applications
\item
  Anderson - Control Theory: A Brief Introduction
\item
  Antsaklis - Linear Systems Primer
\item
  Asadi - Feedback Control Systems MATLAB Simulink Approach
\item
  BURGHES - Control and Optimal Control Theories with Applications
\item
  Babu - Control Systems
\item
  Baillieul - Encyclopedia of Systems and Control (2ed)
\item
  Bavafa-Toosi - Introduction to Linear Control Systems
\item
  Bishop - Teaching Modern Control System Analysis and Design
\item
  Bubnicki - Modern Control Theory
\item
  D'Azzo - Linear Control System Analysis and Design (5ed)
\item
  Diana - Control of Mechanical Systems
\item
  Douglas - Fundamentals of Control Theory
\item
  Doyle - Feedback Control Theory
\item
  Frank - Control Theory Tutorial Basic Concepts
\item
  Franklin, Gene F. - Feedback Control of Dynamic Systems (7ed)
\item
  Gazi - Principles of Signals and Systems
\item
  Ghosh - Control Systems Theory and Applications
\item
  Gupta - Automatic Control Engineering
\item
  Guzman - Automatic Control with Interactive Tools
\item
  HARRIS - Stability of Input-Output Dynamical Systems
\item
  HESPANHA, João P. - Linear Systems Theory
\item
  Haidekker, Mark A. - Linear Feedback Controls: The Essentials (2ed)
\item
  Hallauer, Arthur C. - Linear Time-Invariant Dynamic Systems
\item
  Heaston, Richard - Modern Control Theory
\item
  Häaglund, Tore - Automatic Control: Lecture Notes
\item
  Jagan, S. - Control Systems
\item
  KAMARAJU, K. - Linear Systems (2ed): Analysis and Applications
\item
  Kani, S. - Control System Engineering (second edition)
\item
  Keviczky, László - Control Engineering
\item
  Khalil, Hassan K. - Control Systems: An Introduction
\item
  Koppel, David B. - Introduction to Control Theory
\item
  Krishnaveni, V. - Signals and Systems
\item
  Larminat, Philippe de - Analysis and Control of Linear Systems
\item
  Lathi, B. P. - Linear Systems and Signals (3ed)
\item
  Luna, Maria P. - Advances in Dynamical Systems Theory, Models,
  Algorithms and Applications
\item
  Narasimham, S. - Analysis of Linear Control System
\item
  Ogata, Katsuhiko - Modern Control Engineering (5ed)
\item
  Oppenheim, Alan V. - Signals and Systems (2ed)
\item
  Padmanabhan, A. - Control Systems
\item
  Paraskevopoulos, P. N. - Modern Control Engineering
\item
  Qiu, Li - Introduction to Feedback Control
\item
  Rawlings, James B. - Model Predictive Control (2ed)
\item
  Sivanandam, S. N. - Control Systems Engineering using MATLAB (2)
\item
  Skogestad, Sigurd - Multivariable Feedback Control: Analysis and
  Design
\item
  TRIPATHI, A. - Control System Analysis and Design (2ed)
\item
  Truxal, John G. - Automatic Feedback Control System Synthesis
\item
  Tymerski, Robert - Classical and Modern Control Design with examples
  from power electronics
\item
  Vinagre, Blas M. - Time in Control Theory: On Concepts, Measures and
  Uses
\item
  Williams, Robert L. - Linear State Space Control Systems
\item
  Zabczyk, Jerzy - Mathematical Control Theory: An Introduction (2ed)
\item
  Zerz, Eva - Introduction to Systems and Control Theory
\item
  Zhang, Qingsheng - Quantitative Process Control Theory
\end{itemize}

\end{document}
